% =====================================================================
%  Essay — free LaTeX template by mitaa.dev
%  Compile with pdfLaTeX + Biber (Overleaf does this for you).
%  Locally:  latexmk -pdf main.tex
% =====================================================================
\documentclass[11pt, a4paper]{article}
% =====================================================================
%  Styling of the essay template — mitaa.dev
%  The content lives in main.tex. Change the font, colour and
%  margins here; a regular essay needs none of it.
% =====================================================================

% --- Language and font ----------------------------------------------
\usepackage[T1]{fontenc}
\usepackage[english]{babel}
\usepackage{libertinus}          % Libertinus: serif, sans and maths
\usepackage{microtype}

% --- Page ------------------------------------------------------------
\usepackage[a4paper, margin=2.5cm]{geometry}
\usepackage{parskip}             % space between paragraphs, no indent
\linespread{1.08}

% --- Colour -----------------------------------------------------------
\usepackage{xcolor}
\definecolor{accent}{HTML}{1F5FA8}   % the one accent colour; change to taste
\definecolor{subtle}{HTML}{6B7280}

% --- Content ----------------------------------------------------------
\usepackage{graphicx}
\usepackage{booktabs}
\usepackage{enumitem}
\setlist{itemsep=0.2em, topsep=0.3em}
\usepackage{caption}
\DeclareCaptionFont{accent}{\color{accent}}
\captionsetup{font=small, labelfont={bf,accent}, labelsep=period}

% --- References (APA 7) -------------------------------------------------
\usepackage{csquotes}
\usepackage[style=apa, backend=biber]{biblatex}
\addbibresource{references.bib}

% --- Headings ----------------------------------------------------------
\usepackage{titlesec}
\titleformat{\section}{\sffamily\Large\bfseries\color{accent}}{\thesection}{0.75em}{}
\titleformat{\subsection}{\sffamily\large\bfseries}{\thesubsection}{0.75em}{}
\titlespacing*{\section}{0pt}{1.6em}{0.5em}
\titlespacing*{\subsection}{0pt}{1.2em}{0.3em}

% --- Header and footer ------------------------------------------------
\usepackage{fancyhdr}
\pagestyle{fancy}
\fancyhf{}
\fancyhead[L]{\small\sffamily\color{subtle}\thedoctitle}
\fancyhead[R]{\small\sffamily\color{subtle}\thedocauthor}
\fancyfoot[R]{\small\sffamily\thepage}
\renewcommand{\headrulewidth}{0pt}
\fancypagestyle{plain}{\fancyhf{}\fancyfoot[R]{\small\sffamily\thepage}}

% --- Essay details ---------------------------------------------------
% Fill these in at the top of main.tex; they appear in the title block and header.
% \docfield{name} creates \docname{...} to fill in and \thedocname to print it.
\newcommand{\docfield}[1]{%
  \expandafter\newcommand\csname thedoc#1\endcsname{}%
  \expandafter\newcommand\csname doc#1\endcsname[1]{%
    \expandafter\renewcommand\csname thedoc#1\endcsname{##1}}}
\docfield{title}
\docfield{subtitle}
\docfield{author}
\docfield{studentnumber}
\docfield{course}
\docfield{supervisor}
\docfield{date}
\docfield{words}
\docdate{\today}

% One line of details; empty fields are skipped.
\newcommand{\docmeta}[2]{%
  \ifdefempty{#2}{}{\noindent\makebox[7.5em][l]{\color{subtle}#1}#2\par}}

% The title block at the top of the first page.
\newcommand{\makeheader}{%
  \thispagestyle{plain}%
  \begingroup
    \raggedright
    {\color{accent}\rule{2.5em}{2.5pt}}\par\vspace{1.2em}
    {\sffamily\fontsize{24}{28}\selectfont\bfseries\thedoctitle\par}
    \ifdefempty{\thedocsubtitle}{}{%
      \vspace{0.4em}{\sffamily\large\color{subtle}\thedocsubtitle\par}}
    \vspace{1.2em}
    {\small\sffamily\setlength{\parskip}{0.15em}%
      \docmeta{Author}{\thedocauthor}
      \docmeta{Student number}{\thedocstudentnumber}
      \docmeta{Course}{\thedoccourse}
      \docmeta{Lecturer}{\thedocsupervisor}
      \docmeta{Date}{\thedocdate}
      \docmeta{Words}{\thedocwords}}
    \vspace{1.2em}
    {\color{subtle!40}\rule{\linewidth}{0.4pt}}\par
    \vspace{0.6em}
  \endgroup}

% --- Links (load last) -----------------------------------------
\usepackage[hidelinks]{hyperref}
\usepackage[nameinlink, noabbrev]{cleveref}


% --- Details: fill these in -------------------------------------------
\doctitle{The Title of Your Essay}
\docsubtitle{A subtitle that sums up your thesis in one sentence}
\docauthor{First Last}
\docstudentnumber{1234567}
\doccourse{Course name}
\docsupervisor{Lecturer's name}
\docdate{\today}
\docwords{1,500}               % leave empty if no word count is needed

\begin{document}

\makeheader

% ---------------------------------------------------------------------
\section{Introduction}
% ---------------------------------------------------------------------
Open with a hook that pulls the reader straight into the subject: a
recent event, a striking figure or a familiar situation. Then lead, in a
few sentences, to the question your essay answers.

End the introduction with your \textbf{thesis}: one clear claim that the
rest of the essay argues for. A good thesis is debatable, specific and
can be said in a single sentence.

% ---------------------------------------------------------------------
\section{First argument}
% ---------------------------------------------------------------------
Start every paragraph with a topic sentence that states the argument.
Then support it with evidence: a study, an example or a figure. Cite your
sources in APA style, in parentheses \parencite{apa2020} or in the
running text, as with \textcite{booth2016}.

Put a direct quotation in quotation marks and add the page:
\enquote{Every good research question starts with a problem}
\parencite[p.~51]{lastname2024}. Use quotations sparingly; your own
words show better that you understand the source.

% ---------------------------------------------------------------------
\section{Second argument}
% ---------------------------------------------------------------------
Build on the first argument or show another side of your thesis. A list
can briefly show what is at stake:

\begin{itemize}
  \item a first consideration, short and concrete;
  \item a second consideration that strengthens the argument;
  \item a third consideration that leads to the next part.
\end{itemize}

% ---------------------------------------------------------------------
\section{Counterargument and rebuttal}
% ---------------------------------------------------------------------
A strong essay takes the other side seriously. Describe the best
counterargument as fairly as you can, then show why your thesis still
stands. That convinces more than arguments in favour alone.

% ---------------------------------------------------------------------
\section{Conclusion}
% ---------------------------------------------------------------------
Restate your thesis in different words and sum up the main arguments.
Bring in nothing new, but close with a look ahead: what does your
conclusion mean for the reader, or which question remains open?

\printbibliography[title={References}]

\end{document}
